\documentclass[11pt]{article}

\usepackage{amsmath, amssymb, amsthm}
\usepackage[margin=1in]{geometry}

\DeclareMathOperator{\disc}{disc}

\title{The Sum Discrepancy Algorithm}
\author{}
\date{}

\begin{document}

\maketitle

The ``Sum Discrepancy Algorithm'' is an iterative (or online) approach
to constructing a low-discrepancy coloring for a set system. Its core
idea is to process elements one by one and, at each step, assign a color
that minimizes a specific global objective function, typically the sum
of the current discrepancies of all sets, weighted to prioritize the
most imbalanced ones.

\section{Formal Definition}

Given a finite set system \((X, \mathcal{R})\) where \(|X| = n\) and
\(|\mathcal{R}| = m\). The algorithm processes the elements of \(X\)
in some order \(v_1, v_2, \dots, v_n\). At the \(t\)-th iteration, it
determines the color for \(v_t\) using only the information from the
previously colored elements \(\{v_1, \dots, v_{t-1}\}\).

The decision rule for the color \(\chi(v_t) \in \{-1, +1\}\) is defined
as follows:

\begin{enumerate}
  \item \textbf{Compute Current Discrepancies}: For every set
        \(S \in \mathcal{R}\), calculate its discrepancy
        \(\disc^t(S)\) based on the colors assigned so far.
  \item \textbf{Define a Weight Function}: Apply a weight function to
        these discrepancies. Common choices include:
        \begin{itemize}
          \item \textbf{Power Function}:
                \(\sum_{S \in \mathcal{R}} (\disc^t(S))^c\)
                for some constant \(c\).
          \item \textbf{Exponential Function}:
                \(\sum_{S \in \mathcal{R}} \exp(\eta \cdot \disc^t(S))\)
                for some scaling parameter \(\eta > 0\).
        \end{itemize}
  \item \textbf{Minimize the Objective}: Assign \(v_t\) the color
        (either \(+1\) or \(-1\)) that minimizes the value of the chosen
        weight function after this assignment is hypothetically made.
\end{enumerate}

In essence, the algorithm chooses the color that leads to the smallest
increase in a weighted sum of all set discrepancies, where sets with
larger discrepancies contribute more heavily to the sum.

\section{Purpose}

The primary goal of this algorithm is to provide a constructive method
for finding colorings with low discrepancy. Unlike purely existential
proofs (like Spencer's theorem), this algorithm is designed to be
efficient and online, meaning it can be implemented to produce a
coloring without needing to know the entire set system or the future
elements in advance. By optimizing a global potential function (the sum
of weighted discrepancies) rather than greedily fixing only the most
imbalanced set, it aims to avoid the pitfalls of simpler greedy
strategies and achieve bounds like \(O(\sqrt{n \log m})\).

If you'd like to explore the analysis of such algorithms in more detail,
just let me know.

\end{document}
